\section{Waveform Feature Computation}
\label{sec:ch3-waveform-features}

The arc-detection firmware computes waveform features from current samples
rather than from RMS current alone. This is necessary because a series arc can
leave the load current below a conventional breaker trip level while still
introducing low-current gaps, restrike pulses, irregular zero-crossing
behavior, and mid/high-frequency energy [34], [35], [53], [54]. The firmware
uses a 512-sample frame at a 27kHz target sampling rate, with a 256-sample
hop. A Hann window is applied before the FFT to reduce spectral leakage, which
is standard practice in discrete Fourier analysis [70], [75].

The current RMS value used in the feature frame is computed as
\begin{equation}
  I_{\mathrm{rms}} =
  \sqrt{\frac{1}{N}\sum_{n=0}^{N-1} i[n]^2}.
  \label{eq:feature-rms}
\end{equation}
In (\ref{eq:feature-rms}), $i[n]$ is the calibrated current sample after the
firmware's current conversion path and $N$ is the number of samples in the
frame. RMS current represents the effective load current and is used both as a
protection variable and as a baseline reference for several arc features.

The firmware also computes current THD from the fundamental and harmonic tone
powers:
\begin{equation}
  \mathrm{THD}_I =
  100\sqrt{\frac{\sum_{h=2}^{10}P_h}{P_1}} .
  \label{eq:feature-thd}
\end{equation}
In (\ref{eq:feature-thd}), $P_1$ is the tracked fundamental tone power and
$P_h$ is the tone power of harmonic $h$. The firmware limits the harmonic sum
to available harmonics below Nyquist. THD is useful because arcing and
nonlinear loads can distort the current waveform, although THD alone is not
sufficient for reliable arc detection [35].

The mid/high-frequency spectral flux feature compares the normalized magnitude
spectrum of the present frame with the previous frame:
\begin{equation}
  F_{\mathrm{midhf}} =
  50\sum_{k=k_0}^{k_1}
  \left|
  \frac{M_k}{\sum_{r=k_0}^{k_1}M_r}
  -
  \frac{M_{k,\mathrm{prev}}}{\sum_{r=k_0}^{k_1}M_{r,\mathrm{prev}}}
  \right|.
  \label{eq:feature-spectral-flux}
\end{equation}
In (\ref{eq:feature-spectral-flux}), $M_k$ is the FFT magnitude in the 900Hz
to 9000Hz flux band after skipping bins around mains harmonics. The factor 50
comes from the firmware expression $0.5 \times 100$. This feature increases
when the mid/high-frequency spectrum changes abruptly from one frame to the
next, which is consistent with intermittent arc restrike behavior.

The high-frequency energy delta compares the high-frequency energy share with
its baseline:
\begin{equation}
  \Delta E_{\mathrm{HF}} =
  10\log_{10}
  \left(
  \frac{E_{\mathrm{HF}}/(E_{\mathrm{HF}}+E_{\mathrm{UM}})+\epsilon}
       {B_{\mathrm{HF}}+\epsilon}
  \right).
  \label{eq:feature-hf-delta}
\end{equation}
In (\ref{eq:feature-hf-delta}), $E_{\mathrm{HF}}$ is the 4500--9000Hz band
energy, $E_{\mathrm{UM}}$ is the 900--4500Hz upper-midband energy,
$B_{\mathrm{HF}}$ is the baseline high-frequency share, and $\epsilon$ is a
small numerical offset. This feature detects a rise in high-frequency content
relative to recent non-arcing operation.

The midband residual ratio is computed from FFT energy in the 220--3200Hz
midband after excluding harmonic bins:
\begin{equation}
  R_{\mathrm{mid}} =
  20\log_{10}
  \left(
  \frac{\sqrt{E_{\mathrm{mid,res}}}/N}
       {\max(B_I,\,0.10)}
  \right).
  \label{eq:feature-midband-ratio}
\end{equation}
In (\ref{eq:feature-midband-ratio}), $E_{\mathrm{mid,res}}$ is the residual
midband energy, $N$ is the frame length, and $B_I$ is the baseline RMS current.
The firmware uses a 0.10A floor to prevent division by a near-zero baseline.
This feature emphasizes non-harmonic disturbance relative to the normal current
level.

The low-current ratio measures the portion of the frame close to zero current:
\begin{equation}
  L_{\mathrm{ratio}} =
  \frac{100}{N}\sum_{n=0}^{N-1}
  \mathbf{1}
  \left(
  |i_c[n]| \leq
  \max(0.035,\,0.06\max(I_{\mathrm{rms}},0.10))
  \right).
  \label{eq:feature-low-current-ratio}
\end{equation}
In (\ref{eq:feature-low-current-ratio}), $i_c[n]$ is the cleaned current
sample and $\mathbf{1}(\cdot)$ is an indicator function. The 0.035A floor and
0.06 current fraction were verified from the firmware configuration. A high
low-current ratio can indicate current interruption or gap behavior during a
series arc.

The maximum low-current run is computed as
\begin{equation}
  t_{\mathrm{low,max}} =
  \frac{1000\,N_{\mathrm{run,max}}}{f_s}.
  \label{eq:feature-low-current-run}
\end{equation}
In (\ref{eq:feature-low-current-run}), $N_{\mathrm{run,max}}$ is the longest
consecutive run of samples satisfying the low-current condition in
(\ref{eq:feature-low-current-ratio}) and $f_s$ is the sampling rate in hertz.
The result is expressed in milliseconds and is one of the six deployed
base features in the final arc model.

The RMS-current baseline z-score is
\begin{equation}
  Z_I =
  \frac{|I_{\mathrm{rms}}-\mu_I|}{\sigma_I+10^{-6}} .
  \label{eq:feature-irms-zscore}
\end{equation}
In (\ref{eq:feature-irms-zscore}), $\mu_I$ and $\sigma_I$ are the adaptive
baseline mean and standard deviation maintained by the firmware. The baseline
is frozen when arc-like features or large current steps are detected, which
prevents the baseline from immediately absorbing a fault signature.

The feature groups used by the firmware are summarized in
Table~\ref{tab:ch3-feature-groups}.
\begin{table}[H]
  \centering
  \caption{Waveform features and physical interpretation}
  \label{tab:ch3-feature-groups}
  \begin{tabular}{p{0.26\textwidth}p{0.34\textwidth}p{0.30\textwidth}}
    \toprule
    Feature group & Computation basis & Physical meaning \\
    \midrule
    RMS and z-score & Frame RMS and adaptive baseline & Load magnitude and unusual current change \\
    Low-current gap features & Low-current ratio, zero dwell, maximum low-current run & Current interruption and restrike gap behavior \\
    Spectral features & FFT bands, spectral flux, high-frequency energy delta & Broadband and frame-to-frame disturbance \\
    Residual features & Cleaned waveform minus low-frequency baseline & Non-harmonic spikes and midband disturbance \\
    Shape features & Cycle NMSE, peak fluctuation, half-cycle asymmetry & Irregular cycle shape and asymmetry \\
    Pulse features & Residual pulse count per cycle & Short spikes consistent with arc restrike \\
    \bottomrule
  \end{tabular}
\end{table}

Table~\ref{tab:ch3-feature-groups} shows why the final classifier used a
combination of features. No single feature is expected to identify all arc
faults across appliance families; the classifier instead combines gap,
spectral, residual, and context information.
